%% ============================================================
%% CHAPTER 2 — LITERATURE REVIEW
%% ============================================================

\chapter{Literature Review}
\label{ch:literature}

\section{Overview}
\label{sec:lit_overview}

This chapter surveys three bodies of research that this thesis bridges.
\Cref{sec:lit_integrity} reviews pipeline structural integrity and
fracture/fatigue assessment, tracing the development from linear elastic
fracture mechanics to the two-parameter Failure Assessment Diagram now
codified in BS~7910:2019.  \Cref{sec:lit_games} examines Stackelberg
security game theory and its deployment in critical infrastructure
protection, identifying the absence of physics-grounded payoff functions
as the central methodological gap.  \Cref{sec:lit_adv} covers adversarial
robustness of machine learning systems, with particular attention to the
sparse literature on NDE classifier security.  \Cref{sec:lit_gap}
synthesises all three pillars and locates the specific contribution of
this thesis within the existing scholarly landscape.

%% ------------------------------------------------------------------
\section{Pipeline Structural Integrity and Failure Assessment}
\label{sec:lit_integrity}

\subsection{Fracture Mechanics and Fitness-for-Service}
\label{subsec:lit_fracture}

The theoretical basis for pipeline integrity assessment is \ac{LEFM},
in which a surface crack of half-length $a$ subjected to remote stress
$\sigma$ develops a stress intensity factor $K_\mathrm{I} = Y\sigma\sqrt{\pi a}$,
where $Y$ is a geometry-dependent correction factor~\cite{Anderson2005}.
Fracture occurs when $K_\mathrm{I}$ reaches the material's plane-strain
fracture toughness $K_\mathrm{Ic}$, a single-parameter criterion valid
for small-scale yielding.  For transmission pipelines operating near the
toughness transition temperature of older carbon steels, and for weld
heat-affected zones with reduced toughness, this criterion is conservative
but not always discriminating~\cite{Francis2014}.

The R6 procedure, first published by the Central Electricity Generating
Board in 1976 and continuously revised through the British Energy R6
revision~4 programme~\cite{Anderson2005}, resolved this limitation by
introducing a two-parameter \ac{FAD}: the fracture ratio
$\Kr = K_\mathrm{I}/K_\mathrm{mat}$ tracks proximity to fracture
instability, and the load ratio $\Lr = \sigma_\mathrm{ref}/\sigma_y$
tracks proximity to plastic collapse.  A flaw is deemed acceptable if
and only if the assessment point $(\Lr, \Kr)$ lies within the failure
locus $\Kr \leq f(\Lr)$.  The R6 approach was subsequently harmonised
into three analysis levels of increasing rigour.  The \ac{BS7910}:2019
Level~1 screening uses a conservative linearised locus; Level~2 uses
the Option~1 analytical curve adopted in this thesis; Level~3 employs
material-specific stress--strain data for crack-growth analysis in
high-consequence applications~\cite{BS7910:2019}.

A parallel American standard, \ac{API}~579-1/\ac{ASME}~FFS-1, provides
three broadly equivalent levels~\cite{API579:2021}.  Key differences
relevant to transmission pipelines include: (i)~API~579 incorporates
explicit guidance for locally-thinned areas using the modified FAD
(MFAD) concept; (ii)~BS~7910 provides more extensive Annex~P solutions
for stress intensity factors in cylindrical shells under combined
loading; and (iii)~BS~7910:2019 Annex~B calibrates its partial safety
factors directly against pipeline failure statistics, making it the
preferred standard for probabilistic integrity assessments of the type
conducted in \cref{ch:physics}~\cite{BS7910:2019,API579:2021,Mihm2020}.

The $L_{r,\max}$ cut-off, defined as $(\sigma_y + \sigma_u)/(2\sigma_y)$
in BS~7910, prevents the FAD from being applied in the regime where
strain hardening is exhausted.  For API~5L X60 steel
($\sigma_y = \SI{414}{\mega\pascal}$, $\sigma_u = \SI{517}{\mega\pascal}$),
this yields $L_{r,\max} = 1.15$; for X65
($\sigma_y = \SI{448}{\mega\pascal}$, $\sigma_u = \SI{531}{\mega\pascal}$),
$L_{r,\max} = 1.12$ — values confirmed in the \texttt{fad\_engine.py}
validation suite (\cref{ch:physics}).

\begin{figure}[htbp]
  \centering
  \includegraphics[width=0.78\textwidth]{fig2_fad_grade_comparison}
  \caption{BS~7910:2019 Level~2 Option~1 FAD curves for API~5L grades
    X52 through X70.  Higher yield strength shifts $L_{r,\max}$ leftward,
    reducing the acceptable load ratio range.  Assessment points for
    nominal and critical flaws are shown for X60 (primary grade in the
    Gulf Coast network model).}
  \label{fig:fad_grade_comparison}
\end{figure}

%% -----
\subsection{Fatigue Design of Welded Joints}
\label{subsec:lit_fatigue}

Fatigue failure of welded joints is governed by S-N curves of the form
$N = C \cdot \Delta\sigma^{-m}$, where $N$ is the number of cycles to
failure, $\Delta\sigma$ the stress range, and the constants $C$ and $m$
are weld-class-dependent~\cite{IIW_Fatigue:2016}.  The \ac{IIW}
Recommendations for Fatigue Design of Welded Joints and Components
codify 14 FAT classes — labelled by the characteristic fatigue strength
at $2 \times 10^6$ cycles, in megapascals — from FAT~36 (cover-plated
beams, poorest quality) to FAT~160 (smooth full-penetration welds with
post-weld machining)~\cite{IIW_Fatigue:2016,Murakami2002}.  Pipeline
girth welds qualified to API~1104:2021 acceptance criteria typically
fall in the FAT~71--FAT~90 range, depending on weld quality level and
post-weld heat treatment~\cite{API1104:2021}.

Cumulative fatigue damage is assessed through Miner's linear damage rule:
\begin{equation}
  D = \sum_i \frac{n_i}{N_i(\Delta\sigma_i)} \leq D_\mathrm{cr},
  \label{eq:miners_rule}
\end{equation}
where $n_i$ is the counted number of cycles in stress-range block $i$,
$N_i$ is the corresponding endurance life from the S-N curve, and
$D_\mathrm{cr} = 1.0$ is the critical damage sum in the IIW
formulation (conservatively reduced to $0.5$ for safety-critical
applications under variable amplitude loading)~\cite{IIW_Fatigue:2016}.

The S-N relationship exhibits a bilinear character in the high-cycle
regime: a primary inverse slope $m = 3$ for $N < 10^7$ cycles
transitions to a secondary slope $m = 5$ beyond the constant-amplitude
fatigue limit, following the IIW variable amplitude damage
model~\cite{IIW_Fatigue:2016}.  This bilinear representation reflects
the physical transition from short-crack propagation, dominated by
crack-tip plasticity and Miner accumulation, to the threshold regime
where crack closure begins to retard propagation significantly
\cite{Murakami2002,Anderson2005}.  The \texttt{fatigue\_engine.py}
module implements both slopes with the correct IIW FAT-class lookup
table, validated against published reference endurance values in the
310-test suite.

Post-weld heat treatment (\acs{PWHT}) can improve the FAT class by one
to two levels by relieving residual tensile stresses at the weld toe —
the primary crack initiation site in girth welds under axial cycling.
For in-service transmission pipelines, however, \acs{PWHT} is
operationally impractical for most repair scenarios, making the
as-welded FAT class the default for fleet-level assessments~\cite{API1104:2021,
Francis2014}.  This conservatism is embedded in the calibrated seam-type
hierarchy of \cref{tab:calibrated_scf}: fillet and lap-welded legacy
joints (pre-1970, LF-ERW, associated with poor weld quality and no
\acs{PWHT}) exhibit the highest $P_f$ values in the Monte Carlo engine.

\begin{figure}[htbp]
  \centering
  \includegraphics[width=0.82\textwidth]{fig3_sn_weld_comparison}
  \caption{IIW S-N curves for pipeline-relevant weld classes (FAT~63
    to FAT~100).  A two-order-of-magnitude spread in fatigue life exists
    at $\Delta\sigma = \SI{100}{\mega\pascal}$.  The bilinear transition
    at $N = 10^7$ cycles is evident for all classes.  Weld class
    assignments follow API~1104:2021 acceptance criteria.}
  \label{fig:sn_comparison}
\end{figure}

%% -----
\subsection{Monte Carlo Methods in Reliability Assessment}
\label{subsec:lit_mc}

Probabilistic integrity management applies \ac{MC} simulation to
propagate scatter in material properties, flaw dimensions, and operating
loads into a distribution of safety margins~\cite{Caleyo2002}.  The key
advantage over deterministic assessment is the direct quantification of
failure probability $P_f$, which is comparable against regulatory risk
tolerability criteria (typically $10^{-4}$--$10^{-6}$ per year for
high-consequence areas under \ac{ASME}~B31.8~\cite{ASME_B31_8:2025}).

Caleyo et al.~\cite{Caleyo2002} demonstrated MC integration for pipeline
corrosion growth, sampling from lognormal wall-loss distributions and
propagating to burst failure pressure.  Zhu et al.~\cite{Zhu2013}
extended the framework to strain-based design, coupling MC with
finite-element limit-load solutions for wrinkle defects.  Mo and
Li~\cite{Mo2012} applied MC to fatigue crack growth in girth welds under
variable amplitude pressure cycling, finding that SCF uncertainty
dominates the failure probability variance — a result consistent with the
sensitivity analysis in \cref{fig:pf_scf_sensitivity} of this thesis.
The present work extends this line of research to a multi-segment network
context, where individual segment $P_f$ values computed from
$N = \num{10000}$ MC trials are assembled into a NetworkX graph that
feeds directly into the Stackelberg game payoff calculation in
\cref{sec:physics_game}.

%% ------------------------------------------------------------------
\section{Security Games for Critical Infrastructure}
\label{sec:lit_games}

\subsection{Stackelberg Security Games: Foundations}
\label{subsec:lit_ssg}

A \acf{SCG} models a sequential two-player interaction where a
\emph{defender} (leader) commits to a mixed strategy $\bmc$ first,
after which an \emph{attacker} (follower) best-responds having observed
the marginal coverage probabilities~\cite{Conitzer2006}.  Under the
assumption of attacker rationality, the defender seeks a
\acf{SSE} that maximises her expected utility given the attacker's
rational best-response.  Von Stengel and Zamir established that a
Stackelberg commitment weakly dominates Nash equilibrium for the leader,
providing a rigorous game-theoretic justification for
commitment-based security protocols~\cite{Conitzer2006}.

Paruchuri et al.~\cite{Paruchuri2008} established the
\textsc{DOBSS} LP formulation for Bayesian \acp{SCG} with multiple
attacker types, in which each type $\theta \in \Theta$ is represented
by a prior probability $p_\theta$ and a utility function $\Ua^\theta(i)$.
The Bayesian expected defender utility is the prior-weighted average of
type-conditional utilities at the rational best-response target.
Because the follower best-response is a piecewise function of the
leader's mixed strategy, the problem is non-convex in general;
\textsc{DOBSS} linearises it via an enumeration of candidate optimal
attacker targets, solving one LP per target and selecting the global
optimum~\cite{Paruchuri2008}.  This approach scales polynomially in the
number of targets $n$ and attack types $|\Theta|$, making it tractable
for the 22-target, 3-type instance in \cref{ch:physics}.

Korzhyk et al.~\cite{Korzhyk2011} subsequently showed that under mild
conditions on the utility structure (unique attacker best-response), the
Stackelberg value equals the Nash equilibrium value of the zero-sum
game, providing a mathematical bridge between the two solution concepts.
This result is exploited in \cref{subsec:physics_lp} to establish the
uniqueness of the LP-derived SSE for the physics-informed payoff structure.

\subsection{Deployed Systems}
\label{subsec:lit_deployed}

Real-world deployments have demonstrated that SSG-based patrol scheduling
can meaningfully reduce adversarial incident rates.  ARMOR randomised
K-9 unit and police patrol routes across the four terminals of Los Angeles
International Airport~\cite{An2013}; a post-deployment evaluation by the
\ac{TSA} reported a statistically significant reduction in security
anomalies relative to the human-scheduled baseline.  PROTECT scheduled
U.S.\ Coast Guard patrol vessels along Boston Harbour ferry routes,
incorporating three adversary models — opportunistic, capable, and
sophisticated — analogous to the three attacker types (STRATEGIC,
OPPORTUNISTIC, STATE\_ACTOR) in this thesis~\cite{Tambe2011}.  IRIS
randomised U.S.\ Federal Air Marshal seat assignments on domestic flights,
demonstrating deployment at a national scale with real operational constraints~\cite{Tambe2011}.

A structural limitation shared by all deployed systems is that attacker
payoffs are supplied by subject-matter experts as scalar constants,
without formal derivation from structural or physical models of the
protected assets.  A technically sophisticated adversary who understands
that a fillet-welded legacy joint at maximum operating pressure has
$P_f \approx \num{0.88}$ while a seamless segment has $P_f \approx
\num{0.53}$ — a fourfold difference in consequence-weighted failure
probability per unit of attack effort — will exploit the gap between
the expert-estimated payoffs and physical reality.  The physics-informed
payoff structure developed in \cref{ch:physics} closes this gap directly.

\subsection{Critical Infrastructure and Network Attacks}
\label{subsec:lit_ci}

Network-theoretic analysis of critical infrastructure~\cite{Newman2010,
Albert2000} demonstrates that real-world infrastructure graphs exhibit
scale-free degree distributions, and removal of a small number of
high-betweenness nodes causes disproportionate disruption to global
network connectivity~\cite{Albert2000}.  Brown et al.~\cite{Brown2006}
applied bilevel optimisation to find the worst-case interdiction set for
electric power grids; Lewis~\cite{Lewis2014} extended the framework to
interdependent pipeline and power systems, showing that cascading failures
initiated at pipeline compression stations can propagate to wide-area
power outages.

A key limitation of purely network-theoretic interdiction models is the
assumption of structural equivalence across edges: an attacker maximising
betweenness disruption ignores whether the target pipe segment is likely
to fail under the applied attack mode.  Conversely, a purely structural
approach minimising $P_f$ ignores the network consequence of a given
failure.  The Bayesian SSG payoff structure in this thesis explicitly
combines both dimensions: the STATE\_ACTOR type is parameterised to
weight both $P_f$ and betweenness centrality, while the STRATEGIC type
prioritises asset value and the OPPORTUNISTIC type exploits structural
weakness alone.

%% ------------------------------------------------------------------
\section{Adversarial Machine Learning and NDE Security}
\label{sec:lit_adv}

\subsection{Adversarial Attacks on Neural Networks}
\label{subsec:lit_adv_basics}

Szegedy et al.~\cite{Szegedy2014} first demonstrated that small,
imperceptible perturbations to input features can cause deep neural
networks to misclassify with high confidence.  The underlying mechanism
is the geometry of high-dimensional decision boundaries: in feature
spaces of dimension $d \gg 1$, class boundaries are dense, and the
distance from a correctly classified point to the nearest misclassified
region is $\mathcal{O}(1/\sqrt{d})$ for random inputs, shrinking with
problem dimensionality~\cite{Goodfellow2016}.  Goodfellow et al.\
proposed the \acf{FGSM}~\cite{Goodfellow2015}:
\begin{equation}
  \bm{x}^\mathrm{adv} = \bm{x} + \eps \cdot \mathrm{sign}
    \!\left(\nabla_{\bm{x}} \mathcal{L}(\theta, \bm{x}, y)\right),
  \label{eq:fgsm}
\end{equation}
which adds a single signed gradient step bounded in $L_\infty$ norm by
$\eps$.  The key insight is that the gradient direction, not its
magnitude, determines the adversarial perturbation, making FGSM
computationally efficient for any differentiable loss.

Kurakin et al.~\cite{Kurakin2017} extended FGSM to the iterative
\ac{BIM}, applying $T$ small FGSM steps with $L_\infty$ projection
after each iteration.  BIM generally achieves higher attack success
rates than FGSM at the same $\eps$ budget, though at greater
computational cost.  Madry et al.~\cite{Madry2018} formalised \ac{PGD}
— BIM initialised from a random perturbation within the $L_\infty(\eps)$
ball — as the canonical strongest first-order attack under the $L_\infty$
threat model.  A classifier trained to resist PGD perturbations
(adversarial training) is provably robust to all first-order attacks
within the same threat model~\cite{Madry2018}.  The observation in
\cref{ch:adversarial} that BIM outperforms PGD on the NDE feature set
(ASR $\num{20.1}\% > \num{9.4}\%$) reveals an important domain-specific
deviation from the standard ordering: local gradient curvature, not
random initialisation, dominates in the structured 32-dimensional NDE
feature space.

Carlini and Wagner~\cite{Carlini2017} subsequently showed that many
first-generation defences — input preprocessing, detection networks —
can be circumvented by stronger optimisation-based attacks.  This raises
the bar for what constitutes genuine adversarial robustness in
safety-critical applications, and motivates the quantitative ASR
measurements in \cref{ch:adversarial} rather than qualitative
robustness claims.

\subsection{Machine Learning for NDE}
\label{subsec:lit_nde_ml}

Deep learning methods have been applied to automated weld defect
classification from radiographic~\cite{Munir2019} and ultrasonic
inspection data~\cite{Meyendorf2022}.  Munir et al.~\cite{Munir2019}
achieved $>97\%$ accuracy on a four-class weld defect radiograph dataset
using a ResNet-50 backbone, demonstrating that production-grade
classification accuracy is achievable with moderate training sets.
Meyendorf et al.~\cite{Meyendorf2022} reviewed the state of AI-assisted
NDE, identifying ultrasonic phased-array imaging and time-of-flight
diffraction as the modalities with the most mature ML integration.

Despite this accuracy progress, adversarial robustness of NDE
classifiers has received virtually no treatment in the open literature.
The closest structural analogue is the medical imaging domain, where
adversarial perturbations have been shown to cause misdiagnosis of
malignant lesions as benign — a direct analogue to the CRACK $\to$ CLEAN
misclassification scenario in \cref{ch:adversarial}~\cite{Goodfellow2016}.
The present thesis is, to the author's knowledge, the first to
characterise adversarial attack success rates for an NDE weld defect
classifier and to connect those rates quantitatively to the game-theoretic
consequence of a missed defect through the physics-informed payoff structure.

%% ------------------------------------------------------------------
\section{The Unifying Gap}
\label{sec:lit_gap}

Three decades of structural integrity research have produced rigorous,
standards-backed tools for computing $P_f$ from physical first principles.
Two decades of security game research have produced deployable algorithms
for optimal defender resource allocation.  A decade of adversarial ML
research has characterised the fragility of neural-network inspection
systems.  No published work combines all three into a single, coherent,
end-to-end framework.  \Cref{tab:lit_comparison} summarises the research
landscape and positions this thesis.

\begin{table}[htbp]
  \centering
  \caption{Positioning of this thesis relative to the existing literature.
    A tick (\checkmark) indicates the capability is present;
    a dash (---) indicates it is absent.  No prior work satisfies all
    four requirements simultaneously.}
  \label{tab:lit_comparison}
  \begin{tabular}{@{}lcccc@{}}
    \toprule
    Work & Physics $\Pf$ & SSG payoff & Network & Adversarial NDE \\
    \midrule
    BS~7910 / API~579~\cite{BS7910:2019,API579:2021}
      & \checkmark & --- & --- & --- \\
    ARMOR / PROTECT~\cite{An2013,Tambe2011}
      & --- & \checkmark & --- & --- \\
    Brown et al.~\cite{Brown2006}
      & --- & --- & \checkmark & --- \\
    Munir et al.~\cite{Munir2019}
      & --- & --- & --- & ---\textsuperscript{$a$} \\
    Madry et al.~\cite{Madry2018}
      & --- & --- & --- & \checkmark \\
    \textbf{This thesis}
      & \checkmark & \checkmark & \checkmark & \checkmark \\
    \bottomrule
    \multicolumn{5}{@{}l}{\textsuperscript{$a$}High accuracy achieved; adversarial robustness not assessed.}
  \end{tabular}
\end{table}

The specific gap this thesis fills is twofold.  First, no prior security
game formulation for pipeline infrastructure derives attacker payoffs
from physics-based failure probability: existing models assign payoffs
heuristically, creating exploitable information asymmetry for a
technically sophisticated adversary.  Second, no prior work treats
adversarial manipulation of NDE inspection outputs as a strategic
layer within a game-theoretic resource allocation framework, despite the
direct operational pathway from suppressed defect detection to
undefended physical vulnerability.  Stated formally:

\begin{quote}
  \textit{The absence of a physics-grounded payoff function in Bayesian
  Stackelberg security games for pipeline infrastructure, and the
  consequent absence of any treatment of adversarial NDE manipulation
  as a strategic layer within such games.}
\end{quote}

\noindent The following chapter constructs the remedy, beginning with
the PHMSA empirical calibration that grounds all subsequent quantitative
claims in real fleet incident data.
